\chapter{结果文件与图表清单}

\section{结果文件说明}

项目的核心结果保存在 \texttt{results/} 目录中。正文只引用精简表，完整 CSV 作为电子附录保留。

\begin{itemize}
  \item \texttt{dataset\_compliance\_summary.csv}：数据集合规性摘要。
  \item \texttt{preprocessing\_quality\_check.csv}：预处理质量检查。
  \item \texttt{classification\_cv\_results.csv}：分类调参完整结果。
  \item \texttt{classification\_tuned\_metrics.csv}：分类测试集阈值策略指标。
  \item \texttt{classification\_threshold\_tuning.csv}：验证集阈值选择结果。
  \item \texttt{regression\_target\_comparison.csv}：两个回归目标对比。
  \item \texttt{regression\_digital\_dependence\_metrics.csv}：数字依赖回归指标。
  \item \texttt{regression\_productivity\_metrics.csv}：生产力弱预测指标。
  \item \texttt{clustering\_model\_comparison.csv}：聚类算法和簇数对比。
  \item \texttt{clustering\_lifestyle\_profiles.csv}：完整聚类画像表。
  \item \texttt{pca\_explained\_variance.csv}：PCA 解释方差表。
\end{itemize}

\section{图表文件说明}

主要图表保存在 \texttt{figures/} 目录中，包括 EDA 图、分类评估图、回归评估图、聚类图和 PCA 图。最终报告引用的图片已复制到 \texttt{overleaf\_final/figures/}，用于 Overleaf 自包含编译。

\section{补充图表}

图 \ref{fig:appendix-productivity-prediction} 展示生产力得分真实值与预测值对比，作为弱预测结果的补充证据。图 \ref{fig:appendix-classification-importance} 展示分类任务置换重要性，用于辅助理解分类模型主要依赖的输入变量。

\maybefigure{regression_productivity_observed_vs_predicted.png}{生产力得分真实值与预测值对比}{fig:appendix-productivity-prediction}

\maybefigure{classification_permutation_importance.png}{分类任务置换重要性}{fig:appendix-classification-importance}

\section{PCA 解释方差表}

表 \ref{tab:pca-explained-variance} 列出前 10 个主成分的解释方差比例。PCA 仅用于聚类可视化和辅助理解，不作为分类或回归模型输入。

\begin{table}[H]
\centering
\caption{PCA 解释方差附录表（前 10 个主成分）}
\label{tab:pca-explained-variance}
\begin{tabular}{rrr}
\toprule
主成分 & 解释方差比例 & 累计解释方差 \\
\midrule
PC1 & 0.2694 & 0.2694 \\
PC2 & 0.1547 & 0.4241 \\
PC3 & 0.1366 & 0.5607 \\
PC4 & 0.1258 & 0.6865 \\
PC5 & 0.1030 & 0.7895 \\
PC6 & 0.0759 & 0.8654 \\
PC7 & 0.0430 & 0.9084 \\
PC8 & 0.0425 & 0.9509 \\
PC9 & 0.0274 & 0.9783 \\
PC10 & 0.0152 & 0.9935 \\
\bottomrule
\end{tabular}
\end{table}

